% ============================================================
% NEW CHAPTER: CM/Monster理論と動的宇宙論
% Insert after Part II Chapter 8 (functional equation duality)
% ============================================================

\part*{Part II 補遺: Connes-Marcolli 拡張}
\addcontentsline{toc}{part}{Part II 補遺: Connes-Marcolli 拡張}

\section*{第8.5章\quad Connes-Marcolli 系と Moonshine への道}
\addcontentsline{toc}{section}{第8.5章: Connes-Marcolli系とMoonshineへの道}

Bost-Connes（BC）系は $\mathrm{GL}(1)$（1次元、有理数体 $\Q$）の数論を扱う。
その分配関数は Riemann ゼータ $\zeta(\beta)$ である。
本章では、この枠組みを $\mathrm{GL}(2)$（2次元、楕円曲線とモジュラー形式）に拡張する
\textbf{Connes-Marcolli（CM）系}を導入し、
$j$関数の CM 点での値が物理定数を決定する可能性を探る。

\subsection{BC $\to$ CM: 何が変わるか}

\begin{center}
\begin{tabular}{lll}
\toprule
& BC 系（現行） & CM 系（拡張） \\
\midrule
代数的構造 & GL(1)（可換） & GL(2)（非可換） \\
対称群 & $\Gal(\Q^{\mathrm{ab}}/\Q) \cong \hat{\Z}^*$ &
  $\mathrm{GL}_2(\hat{\Z}) \times \mathrm{GL}_2^+(\Q)$ \\
分配関数 & $\zeta(\beta)$ & $L(f, \beta)$（モジュラー形式の $L$ 関数） \\
$\beta = 1$ での相転移 & $\zeta(1) = \infty$（極） & $L(f, 1)$ 有限（極なし） \\
Moonshine 接続 & なし & $j$ 関数 → Monster 群 \\
\bottomrule
\end{tabular}
\end{center}

\subsection{$j$ 関数と CM 点}

$j(\tau)$ は $\mathrm{SL}(2,\Z) \backslash \mathfrak{H}$ のhauptmodul（主関数）であり、
$q$展開
\begin{equation}
  j(\tau) = q^{-1} + 744 + 196884\,q + 21493760\,q^2 + \cdots
  \quad (q = e^{2\pi i\tau})
\end{equation}
を持つ。複素乗法（CM）点では $j$ は代数的整数をとる：

\begin{center}
\begin{tabular}{crcl}
\toprule
判別式 $D$ & $j$ 値 & $j^{1/3}$ & 体 \\
\midrule
$-3$ & $0$ & $0$ & $\Q(\omega)$（Eisenstein 整数） \\
$-4$ & $1728$ & $12$ & $\Q(i)$（Gauss 整数） \\
$-7$ & $-3375$ & $-15$ & $\Q(\sqrt{-7})$ \\
$-8$ & $8000$ & $20$ & $\Q(\sqrt{-2})$ \\
$-11$ & $-32768$ & $-32$ & $\Q(\sqrt{-11})$ \\
\bottomrule
\end{tabular}
\end{center}

$r_D = j(D)^{1/3}$ を「$j$ 三乗根」と呼ぶ。
$r_{-4} = 12$、$|r_{-7}| = 15$ が物理定数に現れる。

\subsection{3つの結合定数と3つの CM 点}

\begin{keybox}[CM点から結合定数へ]
3つの最小類数1判別式が3つのゲージセクターに対応する：
\begin{align}
  D = -3\;\;(j = 0)&: \quad \text{SU}(3)\;\text{（QCD）} \notag \\
  D = -4\;\;(j = 1728)&: \quad \text{U}(1)\;\text{（電磁気）}
  \quad \Longrightarrow \quad
  \alpha_{\mathrm{EM}} = \frac{4\pi}{j(i)} = \frac{4\pi}{1728} = \frac{1}{137.5}
  \label{eq:alpha_em} \\
  D = -7\;\;(j = -3375)&: \quad \text{電弱混合}
  \quad \Longrightarrow \quad
  \sin^2\theta_W = \frac{3}{8\pi}\left(\frac{5}{4}\right)^{\!3}
  = \frac{375}{512\pi} = 0.233
  \label{eq:sin2tw}
\end{align}
\end{keybox}

\noindent
(\ref{eq:alpha_em}) は観測値 $1/137.036$ と \textbf{0.3\%} 一致。
(\ref{eq:sin2tw}) は観測値 $0.2312$ と \textbf{0.8\%} 一致。

$\sin^2\theta_W$ の公式は次のように分解できる：
\begin{equation}
  \sin^2\theta_W
  = \underbrace{\frac{3}{8}}_{\text{GUT 値 (SU(5))}}
  \times \underbrace{\frac{1}{\pi}\left(\frac{r_{-7}}{r_{-4}}\right)^{\!3}}_{\text{RG running 因子}}
  = \frac{3}{8} \times \frac{(5/4)^3}{\pi}
\end{equation}
$j$ 値の比が GUT スケールから $m_Z$ への\textbf{繰り込み群フロー}を暗号化している。

\subsection{Moonshine: 24 と Monster 群}

$j(\tau) - 744 = \sum c_n q^n$ の係数 $c_n$ は Monster 群 $\mathbb{M}$ の
表現の次元である（Monstrous Moonshine, Borcherds 1992）：
$c_1 = 196884 = 196883 + 1$（$196883 = 47 \times 59 \times 71$、全て Monster 素数）。

\begin{dfnbox}[Moonshine module $V^\natural$]
FLM（1988）構成による頂点作用素代数 $V^\natural$：
\begin{itemize}
\item 中心荷 $c = 24$（= Leech 格子の次元 = $|D_4\text{根}|$ = $\eta^{24} = \Delta$）
\item 自己同型群 $= \mathbb{M}$（Monster 群、$|\mathbb{M}| \approx 8 \times 10^{53}$）
\item 分配関数 $= j(\tau) - 744$
\end{itemize}
\end{dfnbox}

\noindent
K理論で現れた「24」（$|\mathrm{im}(J)_3| = 24$、第7章）は
Moonshine の $c = 24$ と\textbf{同一}であった。
これは BC 系の枠内では説明できなかった構造的起源を与える。

Monster 群の素因数は $\{2, 3, 5, 7, 11, 13, 17, 19, 23, 29, 31, 41, 47, 59, 71\}$
（15個）であり、Euler積増幅に必要な素数集合とほぼ一致する。
Monster は Euler 積の\textbf{自然なカットオフ}を提供する可能性がある。


% ============================================================
\section*{第8.6章\quad モジュラー・クインテッセンス：動的宇宙論}
\addcontentsline{toc}{section}{第8.6章: モジュラー・クインテッセンス}

BC 系では $\Omega_\Lambda = 2\pi/9$ は定数（$w = -1$）であった。
CM/Monster 拡張により、モジュラーパラメータ $\tau$ を
\textbf{動的場}（クインテッセンス場）に昇格させることで、
時間変動するダークエネルギーを記述できる。

\subsection{モジュラーポテンシャル}

Monster CFT $V^\natural$ のパラメータ $\tau$ に対し、
スペクトル作用から得られるポテンシャル $V(\tau)$ は
\textbf{モジュラー不変}（$\mathrm{SL}(2,\Z)$ 不変）でなければならない。
重み0のモジュラー不変関数は $j(\tau)$ のみの関数であるから：
\begin{equation}
  V(\tau) = F\!\left(j(\tau)\right)
  \label{eq:modpot}
\end{equation}
$F$ の形は物理で決まるが、$\tau$ 依存性は $j$ を通じてのみ現れる。

最も単純な選択 $V \propto |j|^{-2/3}$ に対し：

\begin{center}
\begin{tabular}{lrcl}
\toprule
$\tau$ & $j(\tau)$ & $V$ & 宇宙論的時代 \\
\midrule
$i\infty$（カスプ） & $\to \infty$ & $\to 0$ & 初期宇宙（DE なし） \\
$i$（自己双対点） & $1728$ & $V_0$ & \textbf{現在}（$\Omega_\Lambda \approx 0.7$） \\
$\rho = e^{2\pi i/3}$ & $0$ & $\to \infty$ & 遠い未来（de Sitter） \\
\bottomrule
\end{tabular}
\end{center}

\subsection{Friedmann 方程式との結合}

\begin{keyresult}[モジュラー Friedmann 方程式]
$\tau(t)$ を宇宙論的場、$V(\tau)$ をモジュラーポテンシャルとすると：
\begin{align}
  H^2 &= \frac{8\pi G}{3}\left[\rho_m(a) + \frac{1}{2}\dot\tau^2 + V(\tau)\right] \\
  \ddot\tau + 3H\dot\tau + V'(\tau) &= 0
  \label{eq:modfried}
\end{align}
状態方程式：$w_\tau = (\frac{1}{2}\dot\tau^2 - V) / (\frac{1}{2}\dot\tau^2 + V)$。
$\dot\tau = 0$ で $w = -1$（宇宙定数）、$\dot\tau \neq 0$ で $w > -1$（クインテッセンス）。
\end{keyresult}

\subsection{美しい分離}

\begin{keybox}[算術と動力学の分離]
\begin{itemize}
\item \textbf{結合定数}（$\alpha_{\mathrm{EM}}$, $\sin^2\theta_W$）
  $= j$ の \textbf{固定 CM 点}での値 → \textbf{定数}（算術的）
\item \textbf{ダークエネルギー}（$\Omega_\Lambda(z)$, $w(z)$）
  $= j$ の \textbf{モジュラー流}での値 → \textbf{時間進化}（動力学的）
\end{itemize}
同じ $j$ 関数の、\textbf{異なる引数}。
\end{keybox}

\noindent
これにより：
\begin{itemize}
\item $\Delta\alpha/\alpha = 0$ at all $z$（結合定数は赤方偏移によらず不変）
  → クェーサー吸収線スペクトルで検証可能
\item $w(z)$ はモジュラーポテンシャルから決まる（自由パラメータ：$\tau_0$ のみ）
  → DESI Y3 / Euclid で検証可能
\item 宇宙は $\tau \to \rho$（$j \to 0$）で de Sitter 状態に収束
\end{itemize}

\subsection{DESI データとの整合性}

DESI (2024) は $w_0 \approx -0.83$, $w_a \approx -0.75$ を報告し、
ダークエネルギーの時間変動を示唆した。

モジュラー・クインテッセンスでは、$\tau$ が $\rho$（$j = 0$）に向かって
加速するにつれ、ポテンシャルが急峻化し、$w$ が増加する
（より負でなくなる）。これは DESI の $w_0 w_a$ CDM 傾向と
\textbf{定性的に一致}する。

$w_0 w_a$ CDM は2パラメータ（$w_0$, $w_a$）のフィット。
モジュラー・クインテッセンスは1パラメータ（$\tau_0$）で
$w(z)$ 関数全体を予測する → より制約的、より検証可能。

\subsection{宇宙の「地図」としての基本領域}

\begin{verbatim}
       i∞ (cusp)                  ← 初期宇宙 (j→∞, V→0)
        │
        │
  ──────┼──────  Re(τ) = ±1/2
       │ │
      /   \
     /     \                      ← 物質優勢期
    /   i   \
   /    ★    \                    ← 現在 (j=1728, Ω_Λ≈0.7)
  /           \
 /      ρ      \
/       ★       \                 ← 未来 (j→0, de Sitter)
─────────────────
     |τ| = 1
\end{verbatim}

宇宙の歴史は $\mathrm{SL}(2,\Z)$ 基本領域上の\textbf{一本の軌道}。
Monster 群はこの軌道を含む空間の\textbf{対称群}。

\subsection{静的 WB → 動的 WB}

\begin{center}
\begin{tabular}{lll}
\toprule
& 静的 WB（BC 系） & 動的 WB（CM/Monster 系） \\
\midrule
予測数 & 3（$\Omega_\Lambda$, $\alpha_{\mathrm{EM}}$, $\sin^2\theta_W$） &
  $\infty$（$w(z)$, $H(z)$, $\Omega_\Lambda(z)$） \\
時間変動 & なし & $\tau(t)$ のモジュラー流 \\
$w$ & $-1$（定数）or $-0.92$（$\zeta$-quint） &
  $w(z)$ がポテンシャルから決まる \\
$\Delta\alpha/\alpha$ & 未定義 & $= 0$（予測） \\
終状態 & 未定義 & de Sitter（$\tau \to \rho$, $j \to 0$） \\
\bottomrule
\end{tabular}
\end{center}

\newpage
